% !TEX program = xelatex
% AI-effects 프로젝트 작업 노트 --- 대화량 기반 직업별 AI 활용도 지표 2개의 자료, 절차, 주의사항.
% 계획 노트: 실제 자료 수집·분석은 아직 하지 않음. 분모 자료의 공개 수준과 공표 시기는 신청 전 재확인 필요.
% 관련: paper 19(한국 대화량), paper 21(SOC–KSCO 연계표), paper 23(KSCO 결합과 특이사항), report 24(산업별 LCE 자료준비),
%       reference/notes/chang-2018-industry-labor-productivity-convergence-employment-ch3.md(노동투입 측정)
% 컴파일: xelatex (kotex 필요)
\documentclass[11pt]{article}

\usepackage{fontspec}
\usepackage{kotex}
\usepackage[a4paper,margin=2.2cm]{geometry}
\usepackage{longtable}
\usepackage{booktabs}
\usepackage{array}
\usepackage{amsmath}
\usepackage{xcolor}
\usepackage{hyperref}
\usepackage{enumitem}

\hypersetup{colorlinks=true,linkcolor=blue!50!black,urlcolor=blue!50!black}
\newcolumntype{L}[1]{>{\raggedright\arraybackslash}p{#1}}
\setlength{\LTleft}{0pt}
\setlength{\LTright}{0pt}
\renewcommand{\arraystretch}{1.15}
\setlength{\tabcolsep}{4pt}

\title{직업별 AI 활용도 지표의 자료와 준비사항\\[2mm]
\large --- 대화량 기반 활용도 2개: 필요한 자료, 절차, 주의사항 ---}
\author{AI-effects 프로젝트 (작업 노트)}
\date{\today}

\begin{document}
\maketitle

\begin{abstract}
\noindent 한국 직업별 업무용 AI 대화량(paper 23)을 이용해, 사전적 지표인 AI 노출도와 구분되는 실현 지표로서 직업별 AI 활용도 두 가지를 만들 계획이다. 활용도 1은 취업자 1인당 업무용 대화 수, 활용도 2는 AI로 줄어든 시간이 총근로시간에서 차지하는 비중이다. 본 노트는 두 지표의 정의, 필요한 자료, 작업 절차, 해석과 측정상의 주의사항, 작업 전에 정할 사항을 정리한 계획 노트다. 실제 자료 수집과 계산은 아직 하지 않았다. 분모 자료(지역별고용조사 등)의 공개 수준과 공표 시기는 작성자의 지식에 근거한 것이므로 신청 전에 MDIS와 통계청 공표 일정에서 다시 확인해야 한다.
\end{abstract}

\tableofcontents

% ============================================================
\section{지표 정의}
% ============================================================

직업 $k$(KSCO 8차 소분류), 월 $t$, 대화당 메시지 수 가정 $K$에 대해

\begin{align}
\text{활용도1}_{k,t} &= \frac{C_{k,t}(K)}{E_k}
  \qquad \text{(취업자 1인당 월 업무용 대화 수)} \\[1mm]
\text{활용도2}_{k,t} &= \frac{S_{k,t}(K)}{H_k},\qquad
  S_{k,t}(K) = \sum_{s} w^{\text{alloc}}_{s,k}\, C_{s,t}(K)\, \frac{\tau_s}{60}
  \qquad \text{(월 근로시간 중 절감 시간 비중)}
\end{align}

\begin{itemize}[leftmargin=*]
  \item $C_{k,t}(K)$: KSCO 소분류 $k$의 월 업무용 대화 수(\texttt{conv\_kK}). $C_{s,t}(K)$는 SOC 2018 세부 직업 $s$의 값.
  \item $\tau_s$: SOC $s$의 대화당 절감 시간(분, \texttt{time\_saved\_min} = 사람 혼자 걸리는 시간 $-$ AI와 함께 걸리는 시간).
  \item $w^{\text{alloc}}_{s,k}$: SOC $\to$ KSCO 배분 가중치(paper 21, SOC 안 합 1).
  \item $E_k$: 직업 $k$의 취업자 수. $H_k = E_k \times h_k \times 4.345$: 월 총근로시간($h_k$ = 주당 평균 근로시간).
\end{itemize}

두 지표 사이에는 다음 관계가 있다.
\[
\text{활용도2}_{k} = \text{활용도1}_{k} \times \frac{\bar\tau_k / 60}{4.345\, h_k},
\qquad \bar\tau_k = \frac{\sum_s w^{\text{alloc}}_{s,k} C_s \tau_s}{\sum_s w^{\text{alloc}}_{s,k} C_s}
\]
즉 두 지표는 기계적으로 상관이 높고, 차이는 대화당 평균 절감 시간 $\bar\tau_k$와 1인당 근로시간 $h_k$의 직업 간 차이에서만 생긴다. 결과에는 이 분해를 함께 보인다.

% ============================================================
\section{필요한 자료}
% ============================================================

\subsection{분자: 보유 자료}

\begin{longtable}{L{5.2cm}L{10.6cm}}
\toprule
파일 & 내용 \\
\midrule
\endhead
\texttt{data/proc/merged/kor\_ksco3\_conv\_ai\_indicators} & 월 $\times$ KSCO 8차 소분류(167개, 군인 4개 포함), 2026-04$\cdot$05. 업무용 대화량 \texttt{conv\_k1p5}$\sim$\texttt{conv\_k10}, 덮개 변수 \texttt{n\_ksco8}, \texttt{n\_ksco8\_conv}, 비율 지표(절감 시간, 노출도) \\
\texttt{kor\_ksco8\_conv\_ai\_indicators} & 같은 내용의 세분류(494개) 판 \\
\texttt{kor\_soc6\_conv\_ai\_indicators} & SOC 2018 세부 직업 판. 활용도 2의 분자(대화량 $\times$ 절감 시간)는 여기서 곱해야 한다 \\
\bottomrule
\end{longtable}

\subsection{분모: 새로 확보할 자료}

\begin{longtable}{L{2.6cm}L{4.6cm}L{8.6cm}}
\toprule
항목 & 자료 & 비고 \\
\midrule
\endhead
취업자 수 $E_k$ & \textbf{지역별고용조사}(통계청, 반기, MDIS) & KSCO 소분류, 자영업자$\cdot$무급가족종사자 포함, 가중치로 집계. 상반기 조사의 기준 시점이 4월이라 대화량(2026-04$\cdot$05)과 시점이 가깝다 \\
 & 경제활동인구조사 & 직업 단위가 거칠어 검증용 \\
근로시간 $h_k$ & 지역별고용조사 주당 취업시간 & 기본. 가구조사라 사업체조사보다 길게 나온다(\S\ref{sec:caution}) \\
 & 고용형태별근로실태조사(고용노동부) & 직종별 근로시간. 임금근로자만 포함. 강건성 검토용 \\
\bottomrule
\end{longtable}

확인할 사항:
\begin{itemize}[leftmargin=*]
  \item \textbf{시점}: 2026년 상반기 지역별고용조사가 공표되었는지. 없으면 2025년 하반기를 쓴다.
  \item \textbf{분류 차수}: 쓸 연도의 직업 코드가 KSCO 7차인지 8차인지. 7차이면 KSCO 7↔8 연계표가 필요하다(report 24에서 미확보로 정리, \texttt{D:/data/AI/data}에 있다고 기록됨).
  \item \textbf{공개 수준}: MDIS 공개용 자료의 직업 코드 자릿수. 소분류가 원격접근으로만 제공되면 중분류로 낮춘다.
\end{itemize}

% ============================================================
\section{절차(do 파일 구성안)}
% ============================================================

\begin{enumerate}[leftmargin=*]
  \item \texttt{01\_import\_kostat\_regional\_employment.do}: 지역별고용조사 원자료 $\to$ \texttt{ksco3}별 가중 취업자 수, 평균 주당 취업시간, 월 총근로시간($E_k h_k \times 4.345$). 표본 수와 상대표준오차도 함께 저장. 출력 폴더는 report 24의 안(\texttt{data/raw/macro/kostat\_lfs\_regional/}, \texttt{data/proc/macro/kostat\_lfs\_regional/})을 따르되 작업 전에 확정한다.
  \item \texttt{03\_merge\_kor\_conv\_ai\_ksco8.do} 수정: SOC 단계에서 $\texttt{saved\_hr\_kK} = \texttt{conv\_kK} \times \texttt{time\_saved\_min}/60$을 만들고 양 변수(\texttt{amt})에 넣어 $w^{\text{alloc}}$으로 세분류$\cdot$소분류에 배분한다. 절감 시간이 있는 대화의 비중(덮개)도 같이 배분한다. 소분류 파일의 대화량과 절감 시간을 곱하면 안 된다(paper 23: KSCO별 최대 +269\% 차이).
  \item \texttt{03\_merge\_kor\_ksco3\_ai\_usage.do}: 대화량 파일과 분모를 \texttt{ksco3}(str3)로 m:1 결합. 월별 활용도 1$\cdot$2와 두 달 평균.
  \item \texttt{04\_analysis\_kor\_ai\_usage\_index.do}: 분포, 상$\cdot$하위 직업, 노출도(Felten AIOE, Eloundou, observed exposure)와의 상관과 순위 비교, 민감도 분석. 출력 \texttt{results/table/25\_*.\{tex,xlsx\}}.
  \item 결과 노트와 \texttt{log.md} 기록(매칭 키: \texttt{ksco3}, 월, 분모 자료의 연도$\cdot$반기).
\end{enumerate}

% ============================================================
\section{주의사항}\label{sec:caution}
% ============================================================

\subsection{해석}
\begin{itemize}[leftmargin=*]
  \item \textbf{AEI의 직업은 사용자의 직업이 아니라 대화 과업이 속한 직업이다.} 활용도 1은 ``직업 $k$ 종사자가 AI를 얼마나 쓰는가''가 아니라 ``직업 $k$의 과업에 대한 AI 사용량을 그 직업 종사자 수로 나눈 값''이다. 학생이나 다른 직업 종사자의 대화도 분자에 들어간다. 지표 이름과 해석에 이 점을 밝힌다.
  \item \textbf{분자 구성 가정}(paper 19의 B1): 한국 8개 제품 업무용 대화의 직업 구성을 Claude.ai 한국 구성으로 대신했다. Claude는 코딩 비중이 높아 소프트웨어 관련 직업의 활용도가 과대 추정될 수 있다.
  \item \textbf{$K$의 역할}: 모든 직업에 같은 $K$를 쓰므로 순위는 $K$와 무관하고 수준만 $1/K$배로 달라진다. 기준 $K$를 정하고, 수준을 인용할 때는 $K$ 범위를 함께 제시한다.
  \item \textbf{절감 시간}: Claude가 대화별로 추정한 값의 평균이며(실측 아님) 범위는 전 세계 Claude.ai다. 사람 혼자 걸리는 시간이 과대 추정되면 활용도 2도 과대 추정된다. 활용도 2는 절감 가능 시간의 상한에 가까운 값으로 해석한다.
\end{itemize}

\subsection{결측과 덮개}
\begin{itemize}[leftmargin=*]
  \item 대화량이 공개되지 않은 소분류 38개(군인 4개 포함), 일부 세분류만 공개된 소분류 34개(2026-05, paper 23). 일부 공개는 하한값이다. 결측을 0으로 두지 않고 덮개 변수를 함께 보고한다.
  \item 절감 시간이 없는 SOC는 활용도 2의 분자에서 빠진다. 절감 시간이 있는 대화 비중을 따로 계산하고, 대체(대분류 평균 등)할지 뺄지 정한다.
  \item 2026-04 농림어업 미공개분(45-XXXX)은 KSCO로 배분되지 않았다.
\end{itemize}

\subsection{분모}
\begin{itemize}[leftmargin=*]
  \item 취업자가 적은 소분류는 표본오차가 커서 비율이 튄다. 최소 취업자 기준을 두거나 중분류 결과를 병행하고, 표본 수와 상대표준오차를 확인한다.
  \item \textbf{근로시간 측정 차이}: 가구조사의 근로시간은 사업체조사보다 길다. 2017년 경제활동인구조사와 사업체노동력조사의 월평균 노동시간 차이는 제조업 22.7시간, 건설업 47.1시간, 도소매$\cdot$음식숙박업 46.5시간이었다(장인성 2018, 표 3-7). 이 차이가 직업마다 다르면 활용도 2의 분모가 직업별로 다르게 과대 추정된다. 기본은 지역별고용조사 시간 하나로 통일하고, 강건성 검토로 같은 보고서 3장 방식(임금근로자는 사업체조사 시간, 비임금근로자는 가구조사의 상대비율로 보정)을 직업 단위에 적용해 본다. 직업 단위 사업체조사 시간은 고용형태별근로실태조사에만 있다.
  \item \textbf{범위}: 분자에 자영업자의 사용도 들어갈 수 있으므로 분모는 임금근로자가 아닌 전체 취업자로 둔다.
  \item \textbf{기간}: 대화량은 두 달뿐이라 계절성을 통제할 수 없다. 두 달 평균을 기본으로 한다.
\end{itemize}

\subsection{검증}
\begin{itemize}[leftmargin=*]
  \item 전국 합(업무용 대화 합 $\div$ 전체 취업자)을 한국은행 가계조사(2025)의 AI 업무 활용률(51.8\%) 등과 비교해 규모를 점검한다.
  \item observed exposure는 같은 AEI 자료에서 나와 활용도와 기계적으로 상관될 수 있다. 노출도와의 비교는 Felten AIOE, Eloundou처럼 독립적인 지표를 중심으로 한다.
\end{itemize}

% ============================================================
\section{작업 전에 정할 사항}
% ============================================================

\begin{enumerate}[leftmargin=*]
  \item 분모 자료와 연도: 지역별고용조사 2026년 상반기(공표 시) 또는 2025년 하반기
  \item 분석 단위: 소분류(167개) 단독 또는 중분류 병행
  \item 기준 $K$
  \item 절감 시간 결측 처리: 제외 또는 대분류 평균 대체
  \item 활용도 2의 근로시간: 가구조사 시간만 쓸지, 사업체조사 기반 보정도 할지
  \item 새 원자료 폴더 이름
\end{enumerate}

\end{document}
